\subsection*{线性网络的三种等效观点}

由线性电阻和独立电源组成的网络满足叠加原理：多个电源共同作用的电流，等于各电源单独作用时
结果的代数和。保留一个理想电压源时，其余理想电压源置零后用短路替代；理想电流源置零后用
断路替代。叠加只用于电压、电流等线性响应，功率含有平方，不能把各分电源产生的功率直接相加。

从任意两个端口向内看，线性有源网络可化为 Thévenin 等效：开路电压 $U_{\rm oc}$ 串联等效电阻
$R_{\rm th}$；也可化为 Norton 等效：短路电流 $I_{\rm sc}$ 并联同一电阻。二者满足
\[
 U_{\rm oc}=I_{\rm sc}R_{\rm th}.
\]
求 $R_{\rm th}$ 时，可把独立源置零后从端口加试探电压，计算 $R=U/I$；若网络含受控源，受控源
必须保留。电压源串联电阻与电流源并联电阻可以互换，因为它们具有相同的端口伏安关系。

\subsection*{对称、等势与网络变换}

对称网络的关键不是“看起来一样”，而是激励和边界条件也在相应交换下不变。若两个节点由对称性
可证明等势，连接它们的电阻中无电流，可以断开；也可以把等势节点合并而不改变其余电路。平衡
电桥条件 $R_1/R_2=R_3/R_4$ 正是中间两点等势。若电源接法破坏对称，原本几何对称的节点不一定
等势。

三端星形网络与三角形网络等效，要求任意两端之间测得的电阻都相同。若星形臂为 $R_A,R_B,R_C$，
则对应三角形边
\[
 R_{AB}=R_A+R_B+\frac{R_AR_B}{R_C},
\]
其余两式循环置换。反变换时，每条星形臂等于与该端相邻两条三角形边的乘积除以三边总和。
这个变换能把不能直接串并联的桥式局部改写成可逐级化简的结构。

无穷网络利用自相似建立方程。设从入口看到的等效电阻为 $R_\infty$，去掉最前一级后余下部分若与
原网络完全相同，便可用另一个 $R_\infty$ 替代，得到代数方程。方程可能有多个根，应保留正值且
满足有限级网络极限的根。更稳妥的检验是计算前若干级 $R_n$，观察递推是否收敛到所选根。

对于正方体、正四面体和多边形网络，可先按电源端点把节点分组：同组节点由对称性等势，连接它们
的支路无电流；剩余支路再按串并联化简。若直接给每条支路设未知电流，虽原则上可由 Kirchhoff
方程求解，却会掩盖网络结构并显著增加计算量。

